\documentclass[12pt]{article}
\usepackage[ukrainian]{babel}
\usepackage{fontspec}
% --- НАЛАШТУВАННЯ ШРИФТІВ ---
\setmainfont{Schoolbook-Regular.otf}[
    Path = ../,
    BoldFont = Schoolbook-Bold.otf,
    ItalicFont = Schoolbook-Italic.otf,
    BoldItalicFont = Schoolbook-BoldItalic.otf
]
\usepackage[italic]{mathastext}
\usepackage[a4paper,margin=1.8cm,bottom=2cm,top=2cm]{geometry}
\usepackage{amsmath,amssymb}
\usepackage{tikz}
\usepackage{pgfplots}
\pgfplotsset{compat=1.16}
\usetikzlibrary{calc,patterns,angles,quotes,intersections,3d}
\usepackage{xcolor,array,fancyhdr,enumitem,multicol}
\usepackage{colortbl}
\usepackage{tcolorbox}
\tcbuselibrary{skins,breakable}
\usepackage[hidelinks]{hyperref}
\usepackage[firstpage=false, text=@pvtr2525, scale=0.6, color=gray!12]{draftwatermark}

% Заборона розриву інлайн-формул
\binoppenalty=10000
\relpenalty=10000

\definecolor{mainGreen}{RGB}{34, 120, 64}
\definecolor{yearOrange}{RGB}{220, 100, 30}
\definecolor{titleBg}{RGB}{225, 240, 220}
\definecolor{instrBg}{RGB}{255, 240, 220}
\definecolor{instrBorder}{RGB}{220, 150, 80}

% Кольори для діаграми
\definecolor{bar1}{RGB}{200, 120, 20}
\definecolor{bar2}{RGB}{255, 220, 0}
\definecolor{bar3}{RGB}{220, 30, 30}
\definecolor{bar4}{RGB}{128, 0, 128}
\definecolor{bar5}{RGB}{0, 160, 220}
\definecolor{bar6}{RGB}{120, 180, 50}

\pagestyle{fancy}
\fancyhf{}
\renewcommand{\headrulewidth}{0.4pt}
\renewcommand{\footrulewidth}{0.4pt}
\fancyhead[C]{\small\color{gray!80} Директорські тести \textendash{} 16 червня 2026}
\fancyhead[R]{\small\color{mainGreen}\textbf{\href{https://t.me/pvtr2525}{@pvtr2525}}}
\fancyfoot[L]{\small\color{gray!80}\href{https://t.me/pvtr2525}{ТГ: @pvtr2525}}
\fancyfoot[C]{\small\color{gray!80}\thepage}
\fancyfoot[R]{\small\color{gray!80}Директорські тести}

\newcommand{\answerTable}[5]{%
\par\vspace{0.15cm}
\noindent
\begin{tabular}{|*{5}{>{\centering\arraybackslash}m{3cm}|}}
\hline
\rule[-0.15cm]{0pt}{0.6cm}\textbf{А} & \rule[-0.15cm]{0pt}{0.6cm}\textbf{Б} & \rule[-0.15cm]{0pt}{0.6cm}\textbf{В} & \rule[-0.15cm]{0pt}{0.6cm}\textbf{Г} & \rule[-0.15cm]{0pt}{0.6cm}\textbf{Д} \\
\hline
\rule[-0.2cm]{0pt}{0.75cm}#1 & \rule[-0.2cm]{0pt}{0.75cm}#2 & \rule[-0.2cm]{0pt}{0.75cm}#3 & \rule[-0.2cm]{0pt}{0.75cm}#4 & \rule[-0.2cm]{0pt}{0.75cm}#5 \\
\hline
\end{tabular}
\par\vspace{0.25cm}
}

\newcommand{\answerTableTall}[5]{%
\par\vspace{0.15cm}
\noindent
\begin{tabular}{|*{5}{>{\centering\arraybackslash}m{3cm}|}}
\hline
\rule[-0.2cm]{0pt}{0.7cm}\textbf{А} & \rule[-0.2cm]{0pt}{0.7cm}\textbf{Б} & \rule[-0.2cm]{0pt}{0.7cm}\textbf{В} & \rule[-0.2cm]{0pt}{0.7cm}\textbf{Г} & \rule[-0.2cm]{0pt}{0.7cm}\textbf{Д} \\
\hline
\rule[-0.45cm]{0pt}{1.2cm}#1 & \rule[-0.45cm]{0pt}{1.2cm}#2 & \rule[-0.45cm]{0pt}{1.2cm}#3 & \rule[-0.45cm]{0pt}{1.2cm}#4 & \rule[-0.45cm]{0pt}{1.2cm}#5 \\
\hline
\end{tabular}
\par\vspace{0.25cm}
}

\newcommand{\instructionBox}[1]{%
\par\vspace{0.3cm}
\begin{tcolorbox}[colback=instrBg, colframe=instrBorder, boxrule=0.6pt, arc=2pt,
    left=8pt, right=8pt, top=4pt, bottom=4pt]
\centering\small\bfseries #1
\end{tcolorbox}
\par\vspace{0.15cm}
}

\newcommand{\sectionTitle}[1]{%
\par\vspace{0.4cm}
\begin{tcolorbox}[colback=titleBg, colframe=titleBg, boxrule=0pt, arc=8pt, halign=center, fontupper=\Large\bfseries\color{mainGreen}]
#1
\end{tcolorbox}\par\vspace{0.3cm}}
\newcommand{\chapterTitle}[1]{\clearpage\sectionTitle{#1}}

\newcommand{\nmtAnswerBox}{%
\par\vspace{0.2cm}\noindent
Відповідь:\ \ 
\begin{tabular}{@{}*{4}{|>{\centering\arraybackslash}p{0.8cm}}|@{\hspace{0.1cm}\textbf{,}\hspace{0.1cm}}*{3}{|>{\centering\arraybackslash}p{0.8cm}}|@{}}
\hline
\rule[-0.2cm]{0pt}{0.7cm} & & & & & & \\
\hline
\end{tabular}
\par\vspace{0.3cm}}

\newcommand{\matchingGrid}{%
    \begingroup
    \renewcommand{\arraystretch}{1.15}
    \setlength{\tabcolsep}{4pt}
    \footnotesize
    \begin{tabular}{r|c|c|c|c|c|}
         \multicolumn{1}{c}{} & \multicolumn{1}{c}{\textbf{А}} & \multicolumn{1}{c}{\textbf{Б}} & \multicolumn{1}{c}{\textbf{В}} & \multicolumn{1}{c}{\textbf{Г}} & \multicolumn{1}{c}{\textbf{Д}} \\ \cline{2-6}
         \textbf{1} & \phantom{X} & \phantom{X} & \phantom{X} & \phantom{X} & \phantom{X} \\ \cline{2-6}
         \textbf{2} & & & & & \\ \cline{2-6}
         \textbf{3} & & & & & \\ \cline{2-6}
    \end{tabular}
    \endgroup
}

\newcounter{zad}
\newcommand{\zadtask}[1]{\stepcounter{zad}\par\vspace{0.3cm}\noindent\textbf{\thezad.}\ \ #1\par\vspace{0.2cm}}
\newcommand{\zadnum}{\stepcounter{zad}\noindent\textbf{\thezad.}\ \ }

\begin{document}
\thispagestyle{empty}
\vspace*{1.2cm}
\begin{center}
{\Large\bfseries\color{mainGreen} ДИРЕКТОРСЬКІ ТЕСТИ З МАТЕМАТИКИ}\\[0.4cm]
{\Huge\bfseries\color{mainGreen} ВАРІАНТ 16 ЧЕРВНЯ 2026 РОКУ}\\[0.6cm]
{\large\color{mainGreen}\textbf{\href{https://t.me/pvtr2525}{Телеграм: @pvtr2525}}}
\end{center}
\vspace{1cm}

\instructionBox{Завдання 1--15 мають по п'ять варіантів відповіді, з яких лише один правильний. Виберіть правильний варіант відповіді й позначте його.}

% ===== №1 =====
\zadtask{У їдальні продається обід з першою стравою вартістю $92{,}4$~грн. Яку найбільшу кількість таких обідів можна придбати на суму $600$~грн?}
\answerTable{$5$}{$7$}{$4$}{$6$}{$8$}

% ===== №2 =====
\noindent
\begin{minipage}[c]{0.42\textwidth}
\zadnum На діаграмі відображено кількість відвідувачів музею протягом шести робочих днів тижня (з вівторка по неділю). У який із днів тижня кількість відвідувачів була вдвічі більшою, ніж попереднього робочого дня?
\end{minipage}%
\hfill
\begin{minipage}[c]{0.56\textwidth}
\centering
\begin{tikzpicture}[scale=0.9, thick, line join=round]
    % Сітка та осі
    \draw[gray!30, thin] (0,0) grid[xstep=1.2, ystep=0.6] (7.2, 6.0);
    \draw[->, thick] (0,0) -- (0, 6.5);
    \draw[->, thick] (0,0) -- (7.5, 0);
    
    % Підписи Y
    \foreach \y/\val in {0/0, 0.6/100, 1.2/200, 1.8/300, 2.4/400, 3.0/500, 3.6/600, 4.2/700, 4.8/800, 5.4/900, 6.0/1000} {
        \node[left, font=\scriptsize] at (0,\y) {\val};
    }
    \node[rotate=90, above, font=\scriptsize, yshift=1cm] at (0,3) {Кількість відвідувачів};
    
    % Стовпчики
    \filldraw[fill=bar1, draw=none] (0.3,0) rectangle (0.9, 0.6);   % 100
    \filldraw[fill=bar2, draw=none] (1.5,0) rectangle (2.1, 1.8);   % 300
    \filldraw[fill=bar3, draw=none] (2.7,0) rectangle (3.3, 1.8);   % 300
    \filldraw[fill=bar4, draw=none] (3.9,0) rectangle (4.5, 3.6);   % 600
    \filldraw[fill=bar5, draw=none] (5.1,0) rectangle (5.7, 5.7);   % 950
    \filldraw[fill=bar6, draw=none] (6.3,0) rectangle (6.9, 4.8);   % 800
    
    % Підписи X
    \node[below, font=\scriptsize] at (0.6,0) {вівторок};
    \node[below, font=\scriptsize] at (1.8,0) {середа};
    \node[below, font=\scriptsize] at (3.0,0) {четвер};
    \node[below, font=\scriptsize] at (4.2,0) {п'ятниця};
    \node[below, font=\scriptsize] at (5.4,0) {субота};
    \node[below, font=\scriptsize] at (6.6,0) {неділя};
\end{tikzpicture}
\end{minipage}
\par\vspace{0.2cm}
\answerTable{середа}{четвер}{п'ятниця}{субота}{неділя}

% ===== Завдання 3 =====
\noindent
\begin{minipage}[c]{0.56\textwidth}
\zadnum Група туристів рухається в напрямку, що утворює з північним напрямком кут $15^\circ$ (див. рисунок). Знайдіть кут $\alpha$ між напрямком руху туристів та західним напрямком.
\end{minipage}%
\hfill
\begin{minipage}[c]{0.42\textwidth}
\centering
\begin{tikzpicture}[scale=0.85, thick, line join=round]
    \draw[<->, >=latex] (-3,0) node[below, text=cyan!80!black] {Захід} -- (3,0) node[below, text=cyan!80!black] {Схід};
    \draw[<->, >=latex] (0,-2.5) node[below, text=cyan!80!black] {Південь} -- (0,3) node[above, text=cyan!80!black] {Північ};
    
    \draw[->, >=latex] (0,0) -- (75:3.2);
    \draw[<-, >=latex, mainGreen] (180:1.8) arc (180:75:1.8) node[midway, above left, text=black] {$\alpha$};
    
    \draw (90:1.2) arc (90:75:1.2);
    \draw (90:1.35) arc (90:75:1.35);
    \node[above right, font=\small] at (75:1.0) {$15^\circ$};
\end{tikzpicture}
\end{minipage}

\par\vspace{0.2cm}
\answerTable{$75^\circ$}{$105^\circ$}{$115^\circ$}{$125^\circ$}{$95^\circ$}
% ===== №4 =====
\zadtask{Розв'яжіть рівняння $\dfrac{1}{2x} = \dfrac{1}{2 - 3x}$.}
\answerTable{$-2$}{$-0{,}4$}{$2{,}5$}{$0{,}4$}{$2$}

% ===== №5 =====
\noindent
\begin{minipage}[c]{0.54\textwidth}
\zadnum На рисунку зображено правильну шестикутну призму $ABCDEFA_1B_1C_1D_1E_1F_1$. Скільки всього є прямих, що містять бічні ребра призми та є паралельними грані $ABB_1A_1$?
\vspace{0.2cm}
\par\noindent\textbf{А}\ \ жодної
\par\noindent\textbf{Б}\ \ лише дві
\par\noindent\textbf{В}\ \ лише чотири
\par\noindent\textbf{Г}\ \ лише п'ять
\par\noindent\textbf{Д}\ \ шість
\end{minipage}%
\hfill
\begin{minipage}[c]{0.44\textwidth}
\centering
\begin{tikzpicture}[scale=0.85, thick, line join=round]
    % нижня основа
    \coordinate (A) at (0,0);
    \coordinate (B) at (-1.1,0.7);
    \coordinate (C) at (0,1.4);
    \coordinate (D) at (2.1,1.4);
    \coordinate (E) at (3.2,0.7);
    \coordinate (F) at (2.1,0);
    % висота
    \def\h{5}
    \coordinate (A1) at (0,\h);
    \coordinate (B1) at (-1.1,\h+0.7);
    \coordinate (C1) at (0,\h+1.4);
    \coordinate (D1) at (2.1,\h+1.4);
    \coordinate (E1) at (3.2,\h+0.7);
    \coordinate (F1) at (2.1,\h);
    % верхня основа (вся суцільна)
    \draw (A1)--(B1)--(C1)--(D1)--(E1)--(F1)--cycle;
    % нижня основа: передні ребра суцільні, задні пунктиром
    \draw (F)--(A)--(B);
    \draw (E)--(F);
    \draw[dashed] (B)--(C)--(D)--(E);
    % бічні ребра
    \draw (A)--(A1);
    \draw (B)--(B1);
    \draw (E)--(E1);
    \draw (F)--(F1);
    \draw[dashed] (C)--(C1);
    \draw[dashed] (D)--(D1);
    % підписи вершин
    \node[below] at (A) {$A$};
    \node[left] at (B) {$B$};
    \node[below right] at (C) {$C$};
    \node[below right] at (D) {$D$};
    \node[right] at (E) {$E$};
    \node[below] at (F) {$F$};
    \node[left] at (A1) {$A_1$};
    \node[left] at (B1) {$B_1$};
    \node[above left] at (C1) {$C_1$};
    \node[above right] at (D1) {$D_1$};
    \node[right] at (E1) {$E_1$};
    \node[right] at (F1) {$F_1$};
\end{tikzpicture}
\end{minipage}
\par\vspace{0.3cm}
% ===== №6 =====
\zadtask{Скоротіть дріб $\dfrac{4b^2 + 20b}{25 + 10b + b^2}$.}
\answerTableTall{$\dfrac{4}{5+b}$}{$\dfrac{4b}{5}$}{$\dfrac{6}{25}$}{$4b(5+b)$}{$\dfrac{4b}{5+b}$}

% ===== №7 =====
\zadtask{Розв'яжіть нерівність $8^{2x+9} > 8^x$.}
\answerTable{$(-\infty; -3)$}{$(-3; +\infty)$}{$(-9; +\infty)$}{$(9; +\infty)$}{$(-\infty; -9)$}

% ===== Завдання 8 =====
\noindent
\begin{minipage}[c]{0.50\textwidth}
\zadnum На рисунку зображено графік функції $y=f(x)$, визначеної на відрізку $[-5;\, 5]$. На якому проміжку ця функція спадає?
\vspace{0.6cm}

\textbf{А}\quad $[-2;\, 2]$\\
\textbf{Б}\quad $[-5;\, -2]$\\
\textbf{В}\quad $[1;\, 6]$\\
\textbf{Г}\quad $[-3;\, 2]$\\
\textbf{Д}\quad $[2;\, 6]$
\end{minipage}%
\hfill
\begin{minipage}[c]{0.48\textwidth}
\centering
\begin{tikzpicture}[scale=0.6, thick]
    \draw[step=1cm, gray!50, thin] (-6,-4) grid (6,4);
    \draw[->, >=latex, very thick] (-6,0) -- (6.5,0) node[below] {$x$};
    \draw[->, >=latex, very thick] (0,-4) -- (0,4.5) node[left] {$y$};
    
    \node[below left] at (0,0) {$0$};
    \node[below] at (-5,0) {$-5$};
    \node[below] at (5,0) {$5$};
    \node[below] at (1,0) {$1$};
    \node[left] at (0,1) {$1$};
    
    % Графік функції
    \draw[line width=1.2pt, smooth, tension=0.7] plot coordinates {(-5,-3) (-4,0) (-2,3) (0,1) (2,-1) (4,1) (5,3)};
    
    \fill (-5,-3) circle (3pt);
    \fill (5,3) circle (3pt);
    
    \node[above right, fill=white, inner sep=1pt] at (1.5, 1.5) {$y=f(x)$};
\end{tikzpicture}
\end{minipage}

\newpage

% ===== №9 =====
\zadtask{Які з наведених тверджень щодо ромба є правильними?
\vspace{0.1cm}
\begin{enumerate}[label=\textbf{\Roman*.}, align=left, leftmargin=2.6em, labelsep=0.5em, itemsep=0.08cm, topsep=0.06cm]
\item Прямі, на яких лежать діагоналі ромба, ділять кути ромба навпіл.
\item Точка перетину діагоналей ромба рівновіддалена від сторін ромба.
\item Сума протилежних кутів ромба дорівнює $180^\circ$.
\end{enumerate}
\vspace{0.1cm}}
\answerTable{лише I}{лише II}{лише I та II}{лише I та III}{I, II та III}

% ===== №10 =====
\zadtask{Послідовність $(a_n)$ задано формулою $n$-го члена $a_n = 2n^2 - 3$. Визначте четвертий член цієї послідовності.}
\answerTable{$5$}{$7$}{$14$}{$2$}{$18$}

% ===== Завдання 11 =====
\noindent
\begin{minipage}[c]{0.50\textwidth}
\zadnum Біля супермаркету в один ряд вирішили встановити чотири контейнери для сортування сміття (див. рисунок). Скількома різними способами можна розставити ці контейнери в ряд, якщо контейнер «Пластик» обов'язково повинен стояти лівіше від усіх інших?
\end{minipage}%
\hfill
\begin{minipage}[c]{0.48\textwidth}
\centering
\begin{tikzpicture}[scale=0.85, thick, line join=round, font=\sffamily]
    % Кольори смуг та кришок
    \definecolor{bin1}{RGB}{255, 255, 255} % Білий (Пластик)
    \definecolor{bin2}{RGB}{250, 230, 100} % Жовтий (Папір)
    \definecolor{bin3}{RGB}{100, 180, 230} % Блакитний (Скло)
    \definecolor{bin4}{RGB}{210, 120, 90}  % Оранжевий (Різне)
    
    \foreach \x/\col/\txt in {0/bin1/ПЛАСТИК, 2.2/bin2/ПАПІР, 4.4/bin3/СКЛО, 6.6/bin4/РІЗНЕ} {
        % Нижня основа (чорна)
        \fill[black!85] (\x, -0.3) rectangle (\x+2, 0);
        
        % Основний корпус (сірий)
        \fill[gray!45] (\x, 0) rectangle (\x+2, 4);
        
        % Перспектива боковини
        \fill[gray!65] (\x+2, -0.3) -- (\x+2.4, 0.1) -- (\x+2.4, 4.4) -- (\x+2, 4) -- cycle;
        \draw[black] (\x+2, -0.3) -- (\x+2.4, 0.1) -- (\x+2.4, 4.4) -- (\x+2, 4);
        
        % Кришка (кольорова)
        \fill[\col] (\x, 4) rectangle (\x+2, 4.5);
        \draw[black] (\x, 4) rectangle (\x+2, 4.5);
        \fill[\col!80!black] (\x, 4.5) -- (\x+0.4, 4.9) -- (\x+2.4, 4.9) -- (\x+2, 4.5) -- cycle;
        \draw[black] (\x, 4.5) -- (\x+0.4, 4.9) -- (\x+2.4, 4.9) -- (\x+2, 4.5) -- cycle;
        
        % Обведення корпусу та основи
        \draw[black] (\x, -0.3) rectangle (\x+2, 0);
        \draw[black] (\x, 0) rectangle (\x+2, 4);
        
        % Отвір для сміття
        \fill[darkgray] (\x+0.3, 3) rectangle (\x+1.7, 3.6);
        \fill[darkgray!60!black] (\x+0.3, 3) -- (\x+0.5, 3.2) -- (\x+1.5, 3.2) -- (\x+1.7, 3) -- cycle;
        
        % Кольорова вертикальна смуга
        \fill[\col] (\x+0.5, 0.3) rectangle (\x+1.5, 2.7);
        \draw[black!50] (\x+0.5, 0.3) rectangle (\x+1.5, 2.7);
        
        % Текст на смузі
        \node[rotate=90, text=black, font=\bfseries] at (\x+1, 1.5) {\txt};
    }
\end{tikzpicture}
\end{minipage}
\par\vspace{0.3cm}

\answerTable{$6$}{$12$}{$48$}{$24$}{$3$}

% ===== №12 =====
\zadtask{У прямокутній системі координат у просторі задано сферу з центром у точці $(6;\, 2;\, 5)$, яка дотикається до координатної площини $xz$. Визначте діаметр цієї сфери.}
\answerTable{$10$}{$4$}{$2$}{$6$}{$12$}

% ===== №13 =====
\zadtask{Для кутів $\alpha$ і $\beta$ виконується умова $\alpha + \beta = 90^\circ$. Визначте $\sin \beta$, якщо $\cos \alpha = \dfrac{1}{4}$.}
\answerTableTall{$\dfrac{1}{4}$}{$\dfrac{\sqrt{15}}{4}$}{$\dfrac{3}{4}$}{$-\dfrac{1}{4}$}{$-\dfrac{\sqrt{15}}{4}$}

% ===== Завдання 14 =====
\noindent
\begin{minipage}[c]{0.56\textwidth}
\zadnum На рисунку зображено рівнобічну трапецію $ABCD$, на основах $BC$ й $AD$ якої вибрано відповідно точки $M$ і $K$ так, що $BK \parallel CD$ і $KM \parallel AB$. $\angle A = 60^\circ$. $KD = 8$~см, $CM = 2$~см. Визначте площу цієї трапеції.
\end{minipage}%
\hfill
\begin{minipage}[c]{0.42\textwidth}
\centering
\begin{tikzpicture}[scale=0.7, thick, line join=round]
    \coordinate (A) at (0,0);
    \coordinate (D) at (8,0);
    \coordinate (B) at (2,3.46);
    \coordinate (C) at (6,3.46);
    \coordinate (K) at (3.46,0);
    \coordinate (M) at (5.3,3.46);
    
    \draw (A) -- (B) -- (C) -- (D) -- cycle;
    \draw (B) -- (K);
    \draw (K) -- (M);
    
    \pic[draw, angle radius=0.6cm, pic text={$60^\circ$}, angle eccentricity=1.6, font=\small] {angle = D--A--B};

    \node[below left] at (A) {$A$};
    \node[above left] at (B) {$B$};
    \node[above right] at (C) {$C$};
    \node[below right] at (D) {$D$};
    \node[below] at (K) {$K$};
    \node[above] at (M) {$M$};
    
    \fill (M) circle (1.5pt);
    \fill (K) circle (1.5pt);
\end{tikzpicture}
\end{minipage}
\par\vspace{0.2cm}
\answerTable{$11\sqrt{3}$~см$^2$}{$33$~см$^2$}{$33\sqrt{3}$~см$^2$}{$66$~см$^2$}{$66\sqrt{3}$~см$^2$}
% ===== №15 =====
\zadtask{Розв'яжіть систему рівнянь 
$\begin{cases} 
2x + y^3 = -15, \\ 
4x - y^3 = 51. 
\end{cases}$ \\
Якщо $(x_0; y_0)$ --- розв'язок системи, то $y_0 =$}
\answerTable{$-3$}{$-27$}{$6$}{$-9$}{$3$}

\newpage

\instructionBox{У завданнях 16--18 до кожного з трьох рядків інформації, позначених цифрами, доберіть один правильний, на Вашу думку, варіант, позначений буквою.}

% ===== №16 =====
\zadtask{Установіть відповідність між виразом (1--3) та його значенням (А--Д).\par
\vspace{0.3cm}
\noindent
\begin{minipage}[t]{0.45\textwidth}
\textit{Вираз}\par\vspace{0.15cm}
\textbf{1} \quad $\log_2 \dfrac{1}{16}$\\[0.3cm]
\textbf{2} \quad $\dfrac{2}{\sqrt{0{,}25}}$\\[0.3cm]
\textbf{3} \quad $(2^{-1})^2$
\end{minipage}%
\hfill
\begin{minipage}[t]{0.3\textwidth}
\textit{Значення виразу}\par\vspace{0.15cm}
\textbf{А} \quad $-4$\\[0.15cm]
\textbf{Б} \quad $-\dfrac{1}{4}$\\[0.15cm]
\textbf{В} \quad $\dfrac{1}{4}$\\[0.15cm]
\textbf{Г} \quad $2$\\[0.15cm]
\textbf{Д} \quad $4$
\end{minipage}%
\hfill
\begin{minipage}[t]{0.2\textwidth}
\vspace{0.4cm}
\begin{flushright}\matchingGrid\end{flushright}
\end{minipage}}

% ===== №17 =====
\zadtask{Доберіть до функції (1--3) властивість її графіка (А--Д).\par
\vspace{0.3cm}
\noindent
\begin{minipage}[t]{0.30\textwidth}
\textit{Функція}\par\vspace{0.2cm}
\textbf{1} \quad $y = \sqrt{x-2} + 1$\\[0.25cm]
\textbf{2} \quad $y = -5^x$\\[0.25cm]
\textbf{3} \quad $y = x^3$
\par\vspace{0.4cm}
\matchingGrid
\end{minipage}%
\hfill
\begin{minipage}[t]{0.66\textwidth}
\textit{Властивість графіка функції}\par\vspace{0.2cm}
\textbf{А} \quad симетричний відносно початку координат\\[0.15cm]
\textbf{Б} \quad симетричний відносно осі $y$\\[0.15cm]
\textbf{В} \quad розміщений лише в ІІІ та IV координатних чвертях\\[0.15cm]
\textbf{Г} \quad розміщений лише в І координатній чверті\\[0.15cm]
\textbf{Д} \quad проходить через точку $(0; 1)$
\end{minipage}}

% ===== Завдання 18 =====
\zadtask{На рисунку зображено прямокутник $ABCD$, у який вписано коло з центром у точці $O$, що дотикається до сторін $BC$ й $AD$, та два півкола, побудованих на сторонах $AB$ і $CD$ як на діаметрах, що мають лише одну спільну точку з колом. $BC = 12$~см. Установіть відповідність між відрізком (1--3) та його довжиною (А--Д).}
\vspace{0.3cm}
\noindent
\begin{minipage}[c]{0.52\textwidth}
\begin{minipage}[t]{0.48\textwidth}
\textit{Відрізок}\par\vspace{0.15cm}
\textbf{1} \quad радіус кола\\[0.25cm]
\textbf{2} \quad $CD$\\[0.25cm]
\textbf{3} \quad $AO$
\end{minipage}
\hfill
\begin{minipage}[t]{0.48\textwidth}
\textit{Довжина відрізка}\par\vspace{0.15cm}
\textbf{А} \quad $3$~см\\[0.15cm]
\textbf{Б} \quad $4$~см\\[0.15cm]
\textbf{В} \quad $3\sqrt{3}$~см\\[0.15cm]
\textbf{Г} \quad $6$~см\\[0.15cm]
\textbf{Д} \quad $3\sqrt{5}$~см
\end{minipage}
\vspace{0.4cm}

\matchingGrid
\end{minipage}%
\hfill
\begin{minipage}[c]{0.45\textwidth}
\centering
\begin{tikzpicture}[scale=0.55, thick, line join=round]
    \coordinate (A) at (0,0);
    \coordinate (B) at (0,6);
    \coordinate (C) at (12,6);
    \coordinate (D) at (12,0);
    
    \draw (A) rectangle (C);
    
    % Півкола
    \draw (0,0) arc (-90:90:3);
    \draw (12,6) arc (90:270:3);
    
    % Коло в центрі
    \draw (6,3) circle (3);
    \fill (6,3) circle (2.5pt) node[right=2pt] {$O$};
    
    \node[below left] at (A) {$A$};
    \node[above left] at (B) {$B$};
    \node[above right] at (C) {$C$};
    \node[below right] at (D) {$D$};
\end{tikzpicture}
\end{minipage}
\newpage

\instructionBox{Розв'яжіть завдання 19--22. Одержані числові відповіді запишіть у спеціально відведеному місці. Відповідь записуйте лише десятковим дробом, урахувавши положення коми. Знак «мінус» записуйте перед першою цифрою числа.}

% ===== №19 =====
\zadtask{Знайдіть різницю між найбільшим і найменшим значеннями функції $f(x) = x^3 - 3x^2 + 20$ на відрізку $[1;\, 5]$.}
\nmtAnswerBox

% ===== №20 =====
\zadtask{Опубліковане в соціальній мережі відео за перші дві доби сумарно набрало $1980$ вподобайок. За першу добу кількість вподобайок була на $20\,\%$ більшою, ніж за другу. Скільки вподобайок набрало відео за першу добу?}
\nmtAnswerBox

% ===== №21 =====
\noindent
\begin{minipage}[c]{0.56\textwidth}
\zadnum На рисунку зображено куб та циліндр. Об'єм циліндра дорівнює $288\pi$~см$^3$, а площа його осьового перерізу дорівнює $72\sqrt{2}$~см$^2$. Діаметр основи циліндра дорівнює діагоналі грані куба. Визначте об'єм (у см$^3$) куба.
\end{minipage}%
\hfill
\begin{minipage}[c]{0.42\textwidth}
\centering
\begin{tikzpicture}[scale=0.65, thick, line join=round]
    % Куб
    \begin{scope}[shift={(0,0)}]
        \coordinate (A) at (0,0);
        \coordinate (B) at (2.5,0);
        \coordinate (C) at (3.5,1.2);
        \coordinate (D) at (1.0,1.2);
        \coordinate (A1) at (0,2.5);
        \coordinate (B1) at (2.5,2.5);
        \coordinate (C1) at (3.5,3.7);
        \coordinate (D1) at (1.0,3.7);
        
        \draw[dashed] (A) -- (D) -- (C);
        \draw[dashed] (D) -- (D1);
        \draw (A) -- (B) -- (C);
        \draw (A) -- (A1);
        \draw (B) -- (B1);
        \draw (C) -- (C1);
        \draw (A1) -- (B1) -- (C1) -- (D1) -- cycle;
    \end{scope}

    % Циліндр
    \begin{scope}[shift={(6.5,0.5)}]
        \draw (0,2.6) ellipse (1.5 and 0.4);
        \draw[dashed] (-1.5,0) arc (180:0:1.5 and 0.4);
        \draw (-1.5,0) arc (180:360:1.5 and 0.4);
        \draw (-1.5,0) -- (-1.5,2.6);
        \draw (1.5,0) -- (1.5,2.6);
    \end{scope}
\end{tikzpicture}
\end{minipage}
\par\vspace{0.2cm}
\nmtAnswerBox

% ===== №22 =====
\zadtask{Знайдіть суму всіх цілих значень $a$, за кожного з яких корінь рівняння
\[
\dfrac{\log_5(4x - 2a + 1) - 2}{\sqrt{x^2 - a^2}} = 0
\]
є цілим числом.}
\nmtAnswerBox
\newpage

% ===== ВІДПОВІДІ =====
\sectionTitle{ПРАВИЛЬНІ ВІДПОВІДІ}

\noindent\textbf{\large Завдання 1--15 (вибір однієї відповіді)}
\par\vspace{0.2cm}
\begin{multicols}{5}
\noindent
\textbf{1.}~В\\
\textbf{2.}~В\\
\textbf{3.}~Б\\
\textbf{4.}~Г\\
\textbf{5.}~В\\
\textbf{6.}~Д\\
\textbf{7.}~Г\\
\textbf{8.}~А\\
\textbf{9.}~В\\
\textbf{10.}~Д\\
\textbf{11.}~А\\
\textbf{12.}~Б\\
\textbf{13.}~А\\
\textbf{14.}~В\\
\textbf{15.}~А
\end{multicols}

\par\vspace{0.3cm}
\noindent\textbf{\large Завдання 16--18 (відповідність)}
\par\vspace{0.2cm}
\noindent
\textbf{16.}\quad 1--А; 2--Д; 3--В\\[0.15cm]
\textbf{17.}\quad 1--Г; 2--В; 3--А\\[0.15cm]
\textbf{18.}\quad 1--А; 2--Г; 3--Д

\par\vspace{0.3cm}
\noindent\textbf{\large Завдання 19--22 (коротка відповідь)}
\par\vspace{0.2cm}
\begin{multicols}{4}
\noindent
\textbf{19.}~$54$\\
\textbf{20.}~$1080$\\
\textbf{21.}~$512$\\
\textbf{22.}~$28$
\end{multicols}

\end{document}